\documentclass[12pt,a4paper]{article}

% ---- Encoding & Font ----
\usepackage{iftex}
\ifPDFTeX
  \usepackage[utf8]{inputenc}
  \usepackage[T1]{fontenc}
\else
  \usepackage{fontspec}  % XeTeX/LuaTeX (tectonic): Unicode nativo (·, —, ¿, ª, §)
\fi
\usepackage[spanish]{babel}
\usepackage{csquotes}

% ---- Page layout ----
\usepackage[top=2.5cm, bottom=2.5cm, left=2.5cm, right=2.5cm]{geometry}
\usepackage{setspace}
\onehalfspacing

% ---- Math ----
\usepackage{amsmath, amssymb, amsthm}
\usepackage{mathtools}
\usepackage{physics}
\usepackage{esint}  % \oiint

% ---- Graphics ----
\usepackage{graphicx}
\usepackage{tikz}
\usepackage{float}
\usepackage{caption}

% ---- Tables ----
\usepackage{array}
\usepackage{booktabs}
\usepackage{tabularx}
\usepackage{colortbl}

% ---- Colours ----
\usepackage{xcolor}
\definecolor{notebg}{HTML}{FFF3CD}
\definecolor{noteborder}{HTML}{FFC107}
\definecolor{boxred}{HTML}{F8D7DA}
\definecolor{boxredborder}{HTML}{F5C6CB}

% ---- Boxes ----
\usepackage{mdframed}
\usepackage{tcolorbox}
\tcbuselibrary{most}

\newtcolorbox{keybox}[1][]{
  colback=blue!5,
  colframe=blue!70!black,
  arc=4pt,
  boxrule=1pt,
  fonttitle=\bfseries,
  title={#1}
}

\newtcolorbox{warnbox}[1][]{
  colback=boxred,
  colframe=boxredborder,
  arc=4pt,
  boxrule=1pt,
  fonttitle=\bfseries,
  title={#1}
}

\newtcolorbox{notebox}[1][]{
  colback=notebg,
  colframe=noteborder,
  arc=4pt,
  boxrule=1pt,
  fonttitle=\bfseries,
  title={#1}
}

% ---- Hyperlinks & TOC ----
\usepackage[hidelinks]{hyperref}
\usepackage{bookmark}

% ---- Enumerate ----
\usepackage{enumitem}

% ---- Miscellaneous ----
\usepackage{icomma}
\usepackage{siunitx}
\usepackage{textcomp}
\usepackage[bottom]{footmisc}

% ---- Title ----
\title{\textbf{Clase 8: Unidades de Potencial, el Electronvolt, \\
        y Potencial de Cargas Puntuales y Distribuciones Continuas}\\[0.3em]
        \large{Volt y electronvolt, potencial de un dipolo,
              y potencial de barra, anillo y disco}}
\author{Transcripto y expandido --- Curso Física III (OpenFING)}
\date{}

\begin{document}

\maketitle
\thispagestyle{empty}
\newpage

\tableofcontents
\newpage

% ===================================================================
\section{Unidades del potencial: el Volt}
% ===================================================================

El potencial y la diferencia de potencial tienen unidades de
\textbf{energía dividida por carga}. En el SI:
\[
[V] = \frac{\text{J}}{\text{C}} \equiv \text{Volt} \;(\text{V}).
\]

% ===================================================================
\section{El electronvolt}
% ===================================================================

Dando vuelta la relación, el Volt permite definir una \textbf{unidad de
energía}: el \textbf{electronvolt} (eV) es la energía que gana una
partícula de carga $e$ al desplazarse entre dos puntos con
$|\Delta V| = 1\,\text{V}$:
\[
1\,\text{eV} = e \cdot (1\,\text{V}) = 1{,}6\times10^{-19}\,\text{J}.
\]

\begin{notebox}[Escalas del electronvolt]
  Unidad muy pequeña, adecuada a la escala atómica y de partículas:
  \begin{itemize}[nosep]
    \item Energía de ligazón más fuerte electrón--núcleo (hidrógeno):
          $\sim 13\,\text{eV}$.
    \item Procesos químicos: fracciones de eV.
    \item Altas energías: múltiplos \textbf{MeV}, \textbf{GeV}. Como $e$
          aparece por todos lados, estas unidades simplifican las
          expresiones.
  \end{itemize}
\end{notebox}

% ===================================================================
\section{Potencial de un conjunto de cargas puntuales}
% ===================================================================

\subsection{Deducción a partir de la energía potencial}

Para cargas $Q_1,Q_2,\dots$ más una carga de prueba $Q_0$ en $\mathbf{r}$
(convención $U=0$ con cargas infinitamente alejadas):
\[
U = \frac{1}{4\pi\varepsilon_0}\sum_{i<j}\frac{Q_i Q_j}{R_{ij}}.
\]

Separando la parte que depende de $Q_0$:
\[
U = \frac{Q_0}{4\pi\varepsilon_0}\sum_{i\geq 1}\frac{Q_i}{R_{0i}}
  + \frac{1}{4\pi\varepsilon_0}\sum_{1\leq i<j}\frac{Q_i Q_j}{R_{ij}}.
\]

El potencial es esa parte por unidad de carga (se cancela $Q_0$):

\begin{keybox}[Potencial de cargas puntuales]
  \[
  \boxed{V(\mathbf{r}) = \frac{1}{4\pi\varepsilon_0}\sum_{i}\frac{Q_i}{R_{0i}}}
  \]
  Independiente de $Q_0$. Nivel de referencia: $V\to 0$ con $Q_0$
  infinitamente alejado ($R_{0i}\to\infty$).
\end{keybox}

\subsection{Ejemplo: cuadrado de cargas}

Cuadrado de lado $A$ con cargas en los vértices; potencial en el centro
$O$. La distancia del centro a cada vértice (media diagonal) es
$d = A/\sqrt{2}$, y como las cuatro son iguales:
\[
V(O) = \frac{1}{4\pi\varepsilon_0}\cdot\frac{\sqrt{2}}{A}\sum_{i=1}^{4} Q_i.
\]

% ===================================================================
\section{Potencial de un dipolo}
% ===================================================================

\begin{notebox}[Por qué el potencial]
  Es una cantidad \textbf{escalar}: evita ángulos, componentes y
  productos vectoriales. Suele ser mucho más fácil de calcular que el
  campo, y el campo se recupera a partir de él (próxima clase).
\end{notebox}

\subsection{Expresión exacta}

Cargas $+Q$ y $-Q$ separadas por $\mathbf{d}$; punto $P$
\textbf{arbitrario} a distancias $r_+$ y $r_-$:
\[
V(P) = \frac{Q}{4\pi\varepsilon_0}\left(\frac{1}{r_+} - \frac{1}{r_-}\right),
\qquad
r_\pm^2 = r^2 \mp \mathbf{r}\cdot\mathbf{d} + \frac{d^2}{4}.
\]

\subsection{Límite de dipolo puntual (Taylor)}

Para $r \gg d$, el término de carga puntual se cancela (carga neta cero);
con $(1+x)^\alpha \approx 1+\alpha x$:
\[
\frac{1}{r_\pm} \approx \frac{1}{r}\left(1 \pm \frac{1}{2}\frac{\mathbf{r}\cdot\mathbf{d}}{r^2}\right).
\]

Restando y con $\mathbf{p} = Q\mathbf{d}$:

\begin{keybox}[Potencial de un dipolo puntual]
  \[
  \boxed{V(P) = \frac{1}{4\pi\varepsilon_0}\frac{\mathbf{r}\cdot\mathbf{p}}{r^3}
        = \frac{1}{4\pi\varepsilon_0}\frac{p\cos\theta}{r^2}}
  \]
  Decae como $1/r^2$ (más rápido que una carga puntual, $1/r$).
\end{keybox}

Con el mismo esfuerzo que el campo (hecho solo sobre el plano bisector),
se obtuvo el potencial en \textbf{cualquier} punto.

% ===================================================================
\section{Potencial de una distribución continua}
% ===================================================================

Dividiendo en elementos $\Delta Q_i = \rho(\mathbf{r}_i)\Delta V_i$
tratados como cargas puntuales, la suma de Riemann tiende a:

\begin{keybox}[Potencial de una distribución continua]
  \[
  \boxed{V(P) = \frac{1}{4\pi\varepsilon_0}\int_{\text{vol}}
         \frac{\rho(\mathbf{r}')}{r}\, dV'}
  \]
  Análogo con $\lambda\,d\ell$ o $\sigma\,dS$. Al ser escalar, no hay que
  proyectar vectores.
\end{keybox}

% ===================================================================
\section{Ejemplo: barra cargada}
% ===================================================================

Barra de largo $L$, densidad lineal uniforme $\lambda$ sobre el eje $z$;
punto $P$ a distancia $y$ perpendicular, en el centro:
\[
V(P) = \frac{1}{4\pi\varepsilon_0}\int_{-L/2}^{L/2}\frac{\lambda\, dz}{\sqrt{y^2+z^2}}
     = \frac{\lambda}{2\pi\varepsilon_0}\int_{0}^{L/2}\frac{dz}{\sqrt{y^2+z^2}}.
\]

\subsection{Integral por sustitución hiperbólica}

Con $z = y\sinh w$: $\sqrt{y^2+z^2} = y\cosh w$ y $dz = y\cosh w\, dw$,
todo se simplifica:
\[
V(P) = \frac{\lambda}{2\pi\varepsilon_0}\operatorname{arcsinh}\!\left(\frac{L}{2y}\right).
\]

Invirtiendo el seno hiperbólico (cuadrática en $u=e^{w}$) se obtiene
$\operatorname{arcsinh}(X) = \ln\!\big(X + \sqrt{X^2+1}\big)$:

\begin{keybox}[Potencial de una barra]
  \[
  \boxed{V(P) = \frac{\lambda}{2\pi\varepsilon_0}
         \ln\!\left(\frac{L}{2y} + \sqrt{\frac{L^2}{4y^2}+1}\right)}
  \]
\end{keybox}

\subsection{Límite de gran distancia}

Para $y \gg L$, con $\ln(1+x)\approx x$:
\[
V(P) \approx \frac{\lambda}{2\pi\varepsilon_0}\cdot\frac{L}{2y}
     = \frac{Q}{4\pi\varepsilon_0\, y}, \qquad Q = \lambda L,
\]
el potencial de una carga puntual, como debía ser.

\begin{notebox}[Ejercicio]
  Analizar el límite opuesto $y \ll L$ (muy cerca de la barra).
\end{notebox}

% ===================================================================
\section{Ejemplo: anillo cargado}
% ===================================================================

Anillo de radio $R$, densidad lineal uniforme $\lambda$; punto $P$ sobre
el eje a distancia $z$. Cada arco tiene $dQ = \lambda R\, d\theta$ a
distancia constante $\sqrt{R^2+z^2}$; como nada depende de $\theta$
($\int_0^{2\pi}d\theta = 2\pi$):

\begin{keybox}[Potencial de un anillo sobre su eje]
  \[
  \boxed{V(P) = \frac{\lambda R}{2\varepsilon_0\sqrt{R^2+z^2}}}
  \]
\end{keybox}

Para $z \gg R$ tiende al potencial de una carga puntual.

% ===================================================================
\section{Ejemplo: disco cargado}
% ===================================================================

Disco de radio $R$, densidad superficial uniforme $\sigma$; punto $P$
sobre el eje a distancia $z$. Se divide en anillos de radio $r$ y ancho
$dr$ ($dQ = \sigma\,2\pi r\, dr$) y se \textbf{suman los potenciales}:
\[
V(P) = \frac{\sigma}{2\varepsilon_0}\int_0^{R}\frac{r\, dr}{\sqrt{z^2+r^2}}.
\]

Con $u = r^2+z^2$:

\begin{keybox}[Potencial de un disco sobre su eje]
  \[
  \boxed{V(P) = \frac{\sigma}{2\varepsilon_0}\left(\sqrt{R^2+z^2} - |z|\right)}
  \]
\end{keybox}

Para $|z| \ll R$: $\;V(P) \approx \dfrac{\sigma}{2\varepsilon_0}(R - |z|)
+ \mathcal{O}(z^2/R)$.

\subsection{Continuidad del potencial y no derivabilidad}

\begin{itemize}[nosep]
  \item El potencial es \textbf{continuo} en $z=0$ (tiende a
        $V_0 = \sigma R/2\varepsilon_0$ por ambos lados).
  \item \textbf{No es derivable}: la dependencia $-|z|$ da un pico en
        $z=0$.
\end{itemize}

\begin{warnbox}[Anticipo: $V$ y el campo]
  La derivada del potencial es (menos) el campo eléctrico. El campo de
  un disco es \textbf{discontinuo} al cruzarlo (apunta hacia arriba de un
  lado y hacia abajo del otro), lo que se corresponde con el salto de la
  derivada de $V$. El potencial es continuo, pero su derivada no.
\end{warnbox}

% ===================================================================
\section{Anuncio: campo desde el potencial}
% ===================================================================

Próxima clase:
\begin{enumerate}[nosep]
  \item Campo eléctrico a partir del potencial (el \textbf{gradiente}).
  \item \textbf{Superficies equipotenciales}.
  \item Comienzo de \textbf{capacitores / condensadores}.
\end{enumerate}

\end{document}
